% !PRINTED 112
\weberpage{112}
%<<CONT

algebra proper, and for that reason it is more satisfactory to answer the \gesperrt{questions of principle} without their aid, as has been done in the preceding and will also be done later. For application and for a more convenient visualization, however, we shall not dispense with this aid.

If, when $b,c$ are real numbers, we put
\begin{equation}
 b=r\cos\varphi,
\qquad c=r\sin\varphi,
\tag{1}
\end{equation}
then, if $r$ is assumed positive, the angle $\varphi$ is thereby determined up to a multiple of $2\pi$. It will be completely determined if we fix an interval of length $2\pi$ in which it is to lie, for example
\[
 -\pi<\varphi\le \pi.
\]
In geometrical interpretation $r,\varphi$ are the polar coordinates of the point whose rectangular coordinates are $b,c$. The absolute value of the complex magnitude $a=b+ci$ is $r$, and we shall call $\varphi$ the \gesperrt{phase} of this complex magnitude.

The product of two complex magnitudes
\begin{equation}
\begin{aligned}
 a&=r(\cos\varphi+i\sin\varphi),\\
 a'&=r'(\cos\varphi'+i\sin\varphi'),
\end{aligned}
\tag{2}
\end{equation}
is obtained in the form
\begin{equation}
 aa'=rr'[\cos(\varphi+\varphi')+i\sin(\varphi+\varphi')],
\tag{3}
\end{equation}
and from this, if $n$ is an arbitrary positive whole number, or even a negative one, one obtains the theorem named after \gesperrt{Moivre},
\begin{equation}
 (\cos\varphi+i\sin\varphi)^n=\cos n\varphi+i\sin n\varphi,
\tag{4}
\end{equation}
which has led to regarding the magnitude $\cos\varphi+i\sin\varphi$ as an exponential function with imaginary exponent, and hence to setting
\[
 \cos\varphi+i\sin\varphi=e^{i\varphi}
\]
Thereby formulas (3) and (4) take the form
\begin{equation}
 e^{i\varphi}e^{i\varphi'}=e^{i(\varphi+\varphi')},\qquad (e^{i\varphi})^n=e^{in\varphi}.
\tag{5}
\end{equation}
If accordingly we put
\begin{equation}
 a=r(\cos\varphi+i\sin\varphi)
\tag{6}
\end{equation}
and by $\sqrt[n]{r}$ understand the positive value existing and completely determined by § 32, IX, whose $n$th power $=r$, then
%CONT>>
